% paper.tex — LaTeX skeleton for the 83-format openXC7 benchmark paper.
%
% This is a SKELETON, not a full conversion of CATALOG_PAPER_DRAFT.md.
% Section headers mirror the markdown structure; body text is future work.
%
% Build:  pdflatex paper && bibtex paper && pdflatex paper && pdflatex paper

\documentclass{article}

\usepackage[utf8]{inputenc}
\usepackage[T1]{fontenc}
\usepackage[hyphens]{url}
\usepackage{hyperref}
\usepackage{graphicx}
\usepackage{booktabs}
\usepackage{amsmath}
\usepackage{amssymb}

\title{83 Number Formats on Open-Source Silicon: \\
       A Reproducible Benchmark of Decode and Compute on openXC7}

\author{Dmitrii Vasilev\\
        \textit{ORCID:} 0009-0008-4294-6159}

\date{Draft v0.1, 2026-07-14}

\begin{document}

\maketitle

\begin{abstract}
We present a reproducible hardware benchmark of 83 numeric formats from 13
families, including IEEE-754, OCP FP8, Microscaling (MX), posit, takum,
decimal, logarithmic, and the phi-derived GoldenFloat family, implemented on a
Xilinx Artix-7 (XC7A200T, ALINX AX7203) using a fully open toolchain (openXC7:
Yosys + nextpnr-xilinx + Project X-Ray). 71 of 83 formats carry at least one
bit-exact silicon decode cell against an independent exact-arithmetic oracle,
reached through four parameterized decode templates (algebraic, table-2\^{}x,
transcendental-exp-via-tables, truncated-multiply). 16 GoldenFloat compute
cells (GF4-GF32, ADD and MUL) verify bit-exactly on the same fabric; GF64
reaches 70.1\% (359/512) bit-exact on silicon, with the residual failures
diagnosed as a timing-closure issue in the 43-bit barrel shifter of the adder
core, not a logic defect. Because Project X-Ray documents the DSP48E1
hard block as only partially reverse-engineered, every design is synthesized
LUT-only (synth\_xilinx -nodsp); we therefore report LUT counts and
place-and-route yields per cell, including the openXC7-specific result that
wide carry-chain multiplies fail routing while wide BRAM tables route
successfully. The contribution is not a new format and claims no superiority
over posit, takum, or microscaling; it is the breadth of formats
proven on vendor-neutral silicon, the open toolchain methodology,
and a per-cell reproducible evidence chain.
\end{abstract}

\section{Introduction}
\textit{(Placeholder --- see CATALOG\_PAPER\_DRAFT.md \S1. To be expanded.)
The 83-format catalog and conformance packs are described in~\cite{vasilev2026catalog}.}

\section{Background --- the number-format landscape}
\textit{(Placeholder --- see \S2.)}
\subsection{IEEE-754 lineage and the minifloat revival}
\subsection{Tapered precision: posit $\rightarrow$ takum $\rightarrow$ tekum}
\subsection{Block-scaled microscaling}
\subsection{First-principles and $\varphi$-derived floats}
\subsection{The 83-format catalog}

\section{Methodology}
\textit{(Placeholder --- see \S3.)}
\subsection{The openXC7 flow}
\subsection{Evidence chain (per Tier-E cell)}
\subsection{Four decode templates}
\subsection{Tier definitions}
\subsection{Accuracy benchmark methodology}

\section{Results}
\textit{(Placeholder --- see \S4.)}
\subsection{Tier-E matrix}
\subsection{Accuracy benchmark}
\subsection{LUT comparison}
\subsection{The openXC7 routing asymmetry}
\subsection{Three case studies where silicon caught what simulation hid}

\section{Related work}
\textit{(Placeholder --- see \S5.)}

\section{Discussion --- limitations of openXC7}
\textit{(Placeholder --- see \S6.)}

\section{Conclusion}
\textit{(Placeholder --- see \S7.)}

\section*{Reproducibility}
\textit{(Placeholder --- pointers to \texttt{format\_benchmark.py},
\texttt{format\_accuracy\_results.csv}, \texttt{lut\_comparison.md}, and
\texttt{formats\_catalog.t27}.)}

\bibliographystyle{plain}
\bibliography{paper}

\end{document}
